\documentclass[UTF8,a4paper,11pt]{ctexart}

\usepackage{doc/references/vggt_vrfm_report}

\hypersetup{
  pdftitle={从相机轨迹多解到修复速度歧义：VGGT 长短上下文修复的理论基础},
  pdfauthor={VGGT Camera Repair Research Notes}
}
\ReportHeaders{VGGT Camera Repair}{Theory Foundation}

\newcommand{\SE}{\mathrm{SE}}
\newcommand{\Sim}{\mathrm{Sim}}
\newcommand{\Sset}{\mathcal S}
\newcommand{\Mspace}{\mathcal M}
\newcommand{\Qspace}{\mathcal Q}
\newcommand{\argminop}{\operatorname*{arg\,min}}
\newcommand{\diam}{\operatorname{diam}}
\newcommand{\rank}{\operatorname{rank}}

\title{\bfseries\color{GlobalBlue}从相机轨迹多解到修复速度歧义\\[0.45em]
\large VGGT 长短上下文修复的理论基础}
\author{Theory Foundation Note}
\date{2026-08-23\\[0.25em]
\texttt{codex/camera\_solution\_space\_docs\_reframe}}

\begin{document}

\maketitle

\begin{statusbox}
\textbf{当前状态：}\texttt{protocol in progress; no scientific conclusion yet}。
本文只建立概念、证据与模型选择边界。前置实验尚未冻结，因此本文不声称 VGGT
已经出现多模态修复速度，也不声称 V-RFM 已经必要。
\end{statusbox}

\begin{abstract}
本文把两个容易混淆的问题分开：其一是在完整观测与评价固定后，去除 gauge 的
相机轨迹可接受集合有何几何和拓扑；其二是在一条 VGGT 长序列预测周围，多个重叠
短窗口是否给出多个有效且不可安全平均的修复速度。前者回答物理可辨识性，后者才是
Variational Rectified Flow Matching（V-RFM）的直接前置问题。我们以 500 帧长序列、
九个长度 100/stride 50 的短窗口为贯穿示例，区分物理多解、gauge、局部弱约束、
优化 basin、算法输出差异和概率/速度多模态；解释普通 MSE 速度回归为何取条件均值，
以及 V-RFM 如何用序列级 latent 表达同一 data--time location 上的多个速度。最后给出
V0 诊断和四类模型决策。严格的 constant-rank、子水平集、紧性与 minimax 讨论保留在
附录。
\end{abstract}

\tableofcontents
\newpage

\section{先说结论}

我们现在面对两个不同问题。

第一个是\textbf{物理解空间问题}：同一组图像和几何约束全部固定以后，是否仍有
多条物理上不同、评价上都成立的相机轨迹？这个问题需要处理 gauge、可辨识性、
子水平集和连通性。

第二个是更贴近当前 VGGT 修复方案的\textbf{修复速度歧义问题}：一条 500 帧序列
先由 VGGT 整体预测，再切成九个长度 100、stride 50 的短窗口。对于相邻窗口共享的
50 帧，左右两个窗口可能给出两份不同的修复建议。我们要问：两份建议是不是都能
改善长序列？如果都有效，它们能不能安全平均？

\begin{keybox}
V-RFM 主要对应第二个问题。只要同一个数据--时间状态上确实存在多个有效、互相冲突、
而且平均后更差的目标速度，就有直接理由研究 latent-conditioned velocity；不要求先
证明世界中存在两个断开的物理真解。
\end{keybox}

但多个短窗口产生多个输出，只说明我们有多个\textbf{候选}。只有通过统一、冻结、
不参与候选生成的评价后，它们才可能成为多个\textbf{有效修复方向}。因此当前关键
步骤是 V0：多短窗修复速度歧义诊断，而不是直接训练 V-RFM。

\section{从 VGGT 的实际场景开始}

\subsection{一条 500 帧轨迹和九个短窗口}

设 VGGT 对完整 500 帧序列给出长序列预测，第 $i$ 帧的 global camera center 为
\begin{equation}
  c_i^G\in\mathbb R^3,\qquad i=0,\ldots,499.
\end{equation}
用长度 100、stride 50 切出九个窗口：$W_0=0{:}99$，$W_1=50{:}149$，依此类推，
$W_8=400{:}499$。相邻窗口总有 50 帧重叠。

两个窗口的局部坐标 gauge 不同，raw camera centers 不能直接相减。先只用预测值把
第 $k$ 个局部结果对齐到 global prediction，得到 $\widetilde c_i^{(k)}$，再定义
\begin{equation}
  d_k(i)=\widetilde c_i^{(k)}-c_i^G.
  \label{eq:local-residual}
\end{equation}
一个 overlap 上的左窗口建议记为 $d_L$，右窗口建议记为 $d_R$。它们是候选修复方向，
还不是两个有效答案。

\subsection{为什么长短帧关系目前很弱}

长窗口和短窗口只共享部分图像，却同时改变了上下文长度、attention 可见的远距离关系、
预测坐标 gauge、累积漂移、边界帧位置、可见性和有效视差。因此 ``100 帧预测'' 与
``500 帧预测'' 不是同一条件下对同一未知量的重复求解。二者不同不能证明固定观测下
存在多个物理解。

但短窗口仍可作为\textbf{候选修复生成器}：不同局部上下文对同一段 global trajectory
提出多个修改方向。V0 要验证的是这些方向的结构，而不是把窗口输出直接升级为物理
真解。

\subsection{两条问题线不能互相替代}

\begin{table}[htbp]
  \centering
  \small
  \begin{tabularx}{\textwidth}{>{\bfseries}p{2.6cm}X X}
    \toprule
    问题线 & 研究对象 & 可以支持的结论 \\
    \midrule
    物理解空间 & 完整条件固定后，去 gauge 的可接受相机轨迹集合
      & 唯一性、弱可辨识方向、道路分量和能垒 \\
    修复速度分布 & 同一 noisy/global state、flow time 与完整 condition 下的 camera residual velocity
      & 确定性融合、selector、连续冗余或速度多模态 \\
    \bottomrule
  \end{tabularx}
  \caption{物理解空间与修复速度分布是两条相关但不能互相替代的问题线。}
\end{table}

物理解可以近似唯一，但算法仍有多条有效修复路径；物理解也可能存在弱方向，而当前
修复目标仍近似单峰。因此不能用任何一条线替另一条线下结论。

\section{六种“看起来不一样”必须分开}

\begin{figure}[htbp]
  \centering
  \begin{tikzpicture}[
    font=\footnotesize,
    card/.style={rounded corners,draw,thick,text width=4.35cm,
      minimum height=2.05cm,align=left,inner sep=7pt},
    flow/.style={-Latex,thick,draw=MeanGray}
  ]
    \draw[GlobalBlue,very thick] (0,0) .. controls (2,0.4) and (4,-0.3) .. (6,0.2);
    \draw[LeftOrange,very thick] (1.7,0.45) .. controls (2.5,0.95) and (3.1,0.85) .. (3.8,0.45);
    \draw[RightGreen,very thick] (2.6,-0.25) .. controls (3.4,-0.8) and (4.1,-0.55) .. (4.8,-0.2);
    \node[text=GlobalBlue,anchor=west] at (6.15,0.2) {global};
    \node[text=LeftOrange,anchor=east] at (1.55,0.55) {left};
    \node[text=RightGreen,anchor=west] at (4.9,-0.25) {right};
    \node[align=center,text=DarkGray] at (3,-1.05)
      {先看到三条预测曲线不同，仍不知道差异属于下面哪一层};

    \node[card,draw=GlobalBlue,fill=PaleBlue] (physical) at (-4.9,-3.0)
      {\textbf{物理多解}\\同一完整观测允许去 gauge 后仍不同的轨迹。};
    \node[card,draw=GlobalBlue,fill=PaleBlue] (gauge) at (0,-3.0)
      {\textbf{Gauge 冗余}\\同一物理轨迹换了全局坐标系。};
    \node[card,draw=GlobalBlue,fill=PaleBlue] (weak) at (4.9,-3.0)
      {\textbf{局部弱约束}\\沿某些方向改一点，评价变化很慢。};

    \node[card,draw=LeftOrange,fill=PaleOrange] (basin) at (-4.9,-5.55)
      {\textbf{优化 basin}\\不同初值或求解器停在不同局部极小。};
    \node[card,draw=LeftOrange,fill=PaleOrange] (output) at (0,-5.55)
      {\textbf{算法输出差异}\\窗口或上下文不同，预测随之变化。};
    \node[card,draw=VRFMPurple,fill=PalePurple] (mode) at (4.9,-5.55)
      {\textbf{概率/速度多模态}\\同一固定条件下有多个有效速度方向。};

    \draw[flow,dashed] (output.west) -- (basin.east);
    \draw[flow,dashed] (output.east) -- (mode.west);
  \end{tikzpicture}
  \caption{六种歧义的证据层级。多个窗口输出首先只属于“算法输出差异”；通过对齐、
  有效性、方向分离和插值检查后，才可能支持修复速度多模态。它仍不自动上升为物理多解。}
  \label{fig:six-ambiguities}
\end{figure}

\begin{longtable}{>{\raggedright\arraybackslash\bfseries}p{2.5cm}
  >{\raggedright\arraybackslash}p{5.1cm}
  >{\raggedright\arraybackslash}p{6.7cm}}
  \toprule
  现象 & VGGT 中的具体例子 & 不能直接推出什么 \\
  \midrule
  \endfirsthead
  \toprule
  现象 & VGGT 中的具体例子 & 不能直接推出什么 \\
  \midrule
  \endhead
  物理多解 & 退化运动或严重遮挡下，两条去 gauge 后仍远离的轨迹都满足独立几何评价
    & 不能由 long/short 输出不同直接推出 \\
  Gauge 冗余 & 整条轨迹一起旋转、平移，单目时还可能整体缩放
    & 不是新物理解，也不是模型 mode \\
  局部弱约束 & 前向小视差时，某段深度/平移方向难以辨认
    & 不等于存在离散分支 \\
  优化 basin & BA 或轨迹优化从不同初值走到不同停点
    & 不等于可行集道路不连通 \\
  算法输出差异 & 100 帧与 500 帧预测，或两个短窗口预测不同
    & 不等于两个输出都有效 \\
  概率/速度多模态 & $d_L,d_R$ 都有效，但中间平均方向明显更差
    & 不等于多个最终物理真解 \\
  \bottomrule
\end{longtable}

最危险的推理，是从算法输出差异直接跳到物理多解或速度多模态。前置实验的作用就是
把候选放进统一评价，再判断它属于 selector、连续冗余还是速度多模态。

\section{物理解空间：重要，但不是 V-RFM 的唯一门票}

\subsection{“固定观测”不只是图片文件名}

完整条件 $C$ 至少包括图像和 frame identity、顺序与长度、所有预处理、内参和畸变、
时间与 rolling-shutter 假设、动态/天空/遮挡 mask、尺度/重力/地图/锚点，以及全部
residual、权重、robust loss、阈值和随机源。少固定一项，研究对象就发生变化。

\subsection{Gauge 是换坐标，不是新世界}

Gauge 是同一物理重建的全局坐标表示差异\citep{hartley2004multiple,triggs2000bundle}。
有 metric depth 或尺度锚点时，尺度通常不可自由变化；纯单目无尺度锚点时，常见 gauge
包含 $\Sim(3)$；若首帧、重力或地图已固定，剩余 gauge 更小。因此群不能机械写死。

\subsection{四个集合与四个描述量}

设轨迹为 $T$，内参为 $K$，场景/nuisance variables 为 $Z$。对冻结能量
$E(C,T,K,Z)$，定义 profile energy
\begin{equation}
  \bar E(C,T,K)=\inf_Z E(C,T,K,Z).
\end{equation}
需要区分 exact fiber、全局最优集 $\argminop\bar E$、相对近优集
$\bar E\le E^*+\epsilon$ 和绝对可接受集 $\bar E\le h$。实验中的“有效候选”通常
属于最后一类，阈值 $h$ 必须预注册。

正容差集合一般含有厚的开邻域，不能简单叫低维流形。正式描述同时报告：商空间直径
$D_h$、道路分量数 $\beta_0$、远候选的 minimax 高度 $\mathcal H_{h,\delta}$ 和局部
弱可辨识维数 $k_{\rm weak}$。详细数学条件见附录。

这条线仍能回答 VGGT 误差来自不可辨识性还是模型不足，并帮助建立独立 scorer 和
阈值。但它是一条长期线，不再是“只有认证断开物理解才允许 latent”的硬门槛。

\section{修复速度歧义：当前最直接的问题}

\subsection{多解可能出现在哪里}

当前方案中最自然的多解位置不是“500 帧最终轨迹有几个世界”，而是：给定同一个
global trajectory state、同一个 flow time 和同一组完整条件证据，目标 camera-center
residual velocity 是否有多个有效方向？

如果模型只看到左窗口或只看到右窗口，两个输出不同只是条件不同。方法设计必须把
global trajectory、全部 local trajectories、coverage、alignment confidence 和 camera
features 一起作为 condition，再研究共同条件下是否需要 latent。短窗口只是 candidate
generator，不能自己宣布 posterior 多模态。

\subsection{为什么两个端点都好还不够}

对两个候选定义
\begin{equation}
  d(\alpha)=(1-\alpha)d_L+\alpha d_R,
  \qquad \alpha\in[0,1].
  \label{eq:residual-interpolation}
\end{equation}
如果所有内部插值都好，两个端点可能只是同一片连续冗余中的不同位置，确定性平均或
加权融合已经足够。如果两个端点都好，但 $\alpha=0.5$ 或一段内部插值明显更差，
“一条确定性平均速度”才可能落入坏区域。

\subsection{四类结果对应四种决策}

\begin{table}[htbp]
  \centering
  \small
  \begin{tabularx}{\textwidth}{>{\raggedright\arraybackslash}p{4.1cm}X X}
    \toprule
    \normalfont\bfseries 标签 & 含义 & 下一步 \\
    \midrule
    \texttt{NOT\_SUPPORTED} & 没有稳定的双端有效、方向分离事件 & 改候选或评价 \\
    \texttt{SELECTOR\_PROBLEM} & 通常只有一个候选有效 & 做 scorer/selector \\
    \texttt{CONTINUOUS\_}\newline\texttt{REDUNDANCY} & 双端和内部平均都有效 & 确定性融合或连续参数化 \\
    \texttt{MULTIMODAL\_}\newline\texttt{VELOCITY\_}\newline\texttt{SUPPORTED}
      & 双端有效且分离，平均/内部显著更差，跨场景稳定
      & 进入 V-RFM 验证 \\
    \bottomrule
  \end{tabularx}
  \caption{V0 的四类结果和对应工程路线。最后一类支持速度 latent，不证明多个物理真解。}
\end{table}

\section{MSE、Diffusion、Rectified Flow 与 V-RFM}

\subsection{MSE 为什么会平均}

若完全相同输入 $u$ 对应多个目标速度 $v$，确定性网络用 MSE 训练的函数空间最优解为
\begin{equation}
  f^*(u)=\mathbb E[v\mid u].
\end{equation}
若两个有效方向相反，均值可能接近零；如果中间方向恰好不好，单一 MSE 输出会形成
错误折中。

\subsection{Diffusion 能处理离散分支}

Diffusion/score model 可以表达多峰甚至看起来分离的分布\citep{ho2020denoising,
song2021score}。“Diffusion 不能处理离散分支”不是一般结论。需要分开表达能力、训练
数据是否覆盖模式，以及我们研究的是最终样本分布还是 data--time space 的速度分布。

V-RFM 与当前问题更直接对齐，是因为我们怀疑同一 data--time location 上存在多个
target velocities，而普通确定性速度场会取平均；并不是因为 Diffusion “做不到”。
同时，有限样本出现两簇既不证明生成分布 support 断开，也不证明物理解空间断开。

\subsection{普通 Rectified Flow 的速度冲突}

Rectified Flow 用耦合样本构造直线路径\citep{liu2023flowstraight}：
\begin{equation}
  x_t=(1-t)x_0+t x_1,
  \qquad v=x_1-x_0.
\end{equation}
当不同 coupling 在同一 $(x_t,t)$ 给出不同目标速度时，确定性
$v_\theta(x_t,t)$ 在平方误差下学习条件均值。Flow Matching 的一般背景见
\citet{lipman2023flow}。

映射到 VGGT 时，$x_1$ 应是通过冻结协议验收的 target camera-center residual，condition
必须包含全部可部署 global/local evidence。否则只是把错误窗口或标签噪声交给生成模型。

\subsection{V-RFM 改了什么}

V-RFM 引入连续 latent $z$\citep{guo2025variational}：
\begin{equation}
  p(v\mid x_t,t,z)=\mathcal N\!\left(v;
  v_\theta(x_t,t,z),I\right).
  \label{eq:vrfm-velocity}
\end{equation}
对 $z$ 积分后，$p(v\mid x_t,t)$ 成为混合分布。同一个 data--time location 可以有
不同 latent-conditioned velocity。训练时 recognition model
$q_\phi(z\mid x_0,x_1,x_t,t)$ 近似 posterior；推理时从 prior $p(z)$ 采样，并在整次
ODE integration 中保持该 $z$。

在 VGGT 修复里我们计划使用\textbf{序列级 latent}：一条 500 帧轨迹只采样一次，
所有帧和 overlap 共享，避免每帧随机切换修复策略。

\subsection{V-RFM 不负责什么}

V-RFM 不会自动知道哪个短窗口正确，也不会把错误候选变成有效候选。它只负责：当
训练目标确有多个有效速度方向时，不必压成一个均值。候选验收和最终样本选择仍需
独立 scorer、几何能量、置信度或下游任务。如果 V0 得到的是
\texttt{SELECTOR\_PROBLEM}，应该先做 selector。

\section{V0：多短窗修复速度歧义诊断}

V0 只回答：多个重叠短窗口是否稳定地产生多个有效、分离、不可安全平均的 camera
repair velocities？

\subsection{冻结对象与最小插值}

必须冻结 scene/frame identity、九个窗口、prediction-only local-to-global alignment、
每个 scene 只拟合一次的 global-to-GT $\Sim(3)$、统一指标、calibration 阈值和
scene-level bootstrap。GT 只用于 privileged offline diagnosis，不能进入候选产生或
prediction-only alignment。

最小插值网格为
\begin{equation}
  \alpha\in\{0,0.25,0.5,0.75,1\}.
\end{equation}
端点回答候选是否各自有效，中间点回答能否安全融合。还需报告方向 cosine、归一化
RMS separation、逐帧一致率、alignment residual、translation error 和 relative
translation error。

\subsection{负控制与 GO 门槛}

负控制至少包含 self-pair、纯 gauge copy、随机错误窗口、残差取反、小扰动和退化
对齐。只有以下条件跨场景稳定，才进入 V-RFM：
\begin{enumerate}
  \item $d_L,d_R$ 对齐可靠；
  \item 两端都通过冻结有效性门槛；
  \item 两方向按预注册 separation 指标显著分开；
  \item 平均或预注册内部插值显著变差；
  \item scene-level bootstrap 不依赖少数特殊场景。
\end{enumerate}
前置实验未结束前，本文不声明这些条件已经满足。

\section{证据等级与安全结论语言}

证据具有不对称性：找到一个远距离有效对，可以反证“所有有效解都近似唯一”；找到
一条完整低能路径，可以证明这对候选在同一道路分量；没找到远解不能证明唯一；路径
搜索失败不能证明断开。对应到 V0，两个模型输出不能证明都有效，两个有效端点也不能
跳过内部插值直接称速度多模态。

\begin{table}[htbp]
  \centering
  \small
  \begin{tabularx}{\textwidth}{p{4.2cm}X X}
    \toprule
    证据 & 可以写 & 不应写 \\
    \midrule
    long/short 输出不同 & 上下文敏感，产生候选修复差异 & 发现多个物理解 \\
    一个候选好、一个差 & 呈现 selector problem & 存在双模态有效速度 \\
    双端和内部都好 & 存在连续可融合冗余 & 平均必然失败 \\
    双端好、内部差且跨场景稳定 & 支持多模态修复速度 & 证明多个物理真轨迹 \\
    路径搜索失败 & 尚未找到低能路径 & 可行集已经断开 \\
    \bottomrule
  \end{tabularx}
  \caption{从观测证据到结论语言的安全边界。}
\end{table}

\section{研究路线}

\begin{center}
\resizebox{0.96\textwidth}{!}{%
\begin{tikzpicture}[
  node distance=7mm,
  block/.style={rounded corners,draw=GlobalBlue,fill=PaleBlue,
    text width=12.5cm,align=center,minimum height=9mm},
  decision/.style={rounded corners,draw=VRFMPurple,fill=PalePurple,
    text width=4.2cm,align=center,minimum height=12mm},
  arrow/.style={-Latex,thick,draw=MeanGray}
]
  \node[block] (s0) {Stage 0：概念、术语与证据边界};
  \node[block,below=of s0] (v0) {V0：多短窗修复速度歧义诊断};
  \node[decision,below left=11mm and -4mm of v0] (det) {NOT SUPPORTED /\\CONTINUOUS\\确定性融合};
  \node[decision,below=11mm of v0] (sel) {SELECTOR\\PROBLEM\\scorer / selector};
  \node[decision,below right=11mm and -4mm of v0] (vrfm) {MULTIMODAL\\VELOCITY\\V-RFM 验证};
  \node[block,below=18mm of sel] (long) {并行长期线：独立能量与合成控制 $\rightarrow$
  局部弱可辨识性 $\rightarrow$ 全局路径与 minimax 能垒 $\rightarrow$ 跨实例 prevalence};
  \draw[arrow] (s0) -- (v0);
  \draw[arrow] (v0) -- (det);
  \draw[arrow] (v0) -- (sel);
  \draw[arrow] (v0) -- (vrfm);
  \draw[arrow,dashed] (v0.south) -- (long.north);
\end{tikzpicture}
}
\end{center}

这条路线把“是否需要 V-RFM”的直接证据提前到 V0，同时保留完整物理解空间研究。
两条线相互校准，但不互相绑架。

\section{当前结论}

\begin{enumerate}
  \item long-vs-short 关系很弱，不能证明固定观测下的物理多解；
  \item 多个 overlap short windows 提供了自然的多候选来源；
  \item 候选是否构成多模态修复速度，取决于双端有效性、方向分离和内部插值；
  \item Diffusion 可以表示多峰分布；V-RFM 的特殊价值是直接建模 data--time space
    的多模态速度；
  \item 平均安全就先做确定性融合；只有一个候选有效就先做 selector；稳定出现
    “端点好、平均差”才进入 V-RFM。
\end{enumerate}

\appendix

\section{状态、观测映射与 profile energy}

设完整状态
\begin{equation}
  x=(T,K,Z)\in\Mspace,
\end{equation}
其中 $T\in\SE(3)^N$ 是相机轨迹，$K$ 是内参，$Z$ 汇总场景几何、深度、对应关系、
动态 mask 等 nuisance variables。理想观测映射为 $F:\Mspace\to\mathcal Y$。
给定观测 $y$，exact fiber 是 $F^{-1}(y)$。若 gauge 群 $G$ 满足
$F(g\cdot x)=F(x)$，物理对象为 $F^{-1}(y)/G$。

冻结能量可写为
\begin{equation}
  E(C,T,K,Z)=\sum_r w_r\rho_r\!\left(e_r(C,T,K,Z)\right)+R(T,K,Z).
\end{equation}
若不同候选使用不同 $Z$ 优化策略，就不再共享一个 profile energy，不能直接比较。

\section{Gauge、商空间与距离}

轨迹等价类记为 $[T]$，商空间距离可写为
\begin{equation}
  d_Q([T_1],[T_2])=\inf_{g\in G}d_{\mathcal T}(T_1,g\cdot T_2).
\end{equation}
$d_{\mathcal T}$ 必须在合法 $\SE(3)$ 表示上定义，并明确旋转/平移权重、时间采样和
尺度处理。四元数 $q$ 与 $-q$、Euler 角周期、raw quaternion 未归一化是表示问题。
VGGT 的 \texttt{pose\_enc\_list} 是同一次确定性 Camera Head 的迭代状态，也不是多个
posterior samples\citep{wang2025vggt}。

\section{Exact fiber 的局部维数}

若 $F$ 在 $x^*$ 附近光滑且 Jacobian 秩局部恒为 $r$，constant-rank theorem
给出\citep{lee2013smooth}
\begin{equation}
  \dim F^{-1}(F(x^*))=\dim\Mspace-r.
\end{equation}
进一步除去局部自由作用的 gauge 维数后，才得到物理 exact fiber 的局部维数。秩必须
在邻域稳定；奇点处 null space 维数可能跳变；一阶 null direction 不保证能积分成
长距离可行曲线。

\section{正容差集合为什么通常是厚的}

令
\begin{equation}
  \Sset_h=\{q\in\Qspace:\bar E(q)\le h\}.
\end{equation}
若 $\bar E$ 连续且存在严格可行点 $q_0$ 满足 $\bar E(q_0)<h$，则 $q_0$ 附近存在
开邻域仍满足 $\bar E<h$。因此 $\Sset_h$ 局部通常是满维厚集合，而不是低维流形。

若 gauge 固定后的物理域闭，$\bar E$ 下半连续且子水平集有界，则 $\Sset_h$ 紧；连续
能量在紧集上取得最小值。这些条件需要逐项验证，不能因为优化器停止就声称找到全局
最优。

\section{直径、道路分量、minimax 与弱方向}

商空间直径为
\begin{equation}
  D_h=\sup_{q_1,q_2\in\Sset_h}d_Q(q_1,q_2).
\end{equation}
大 $D_h$ 反证近似唯一，但不说明集合连通。道路连通分量数 $\beta_0(\Sset_h)$ 不能
由有限采样的聚类数替代\citep{edelsbrunner2010computational}。

两候选间的 minimax 高度定义为
\begin{equation}
  H(q_a,q_b)=\inf_{\gamma(0)=q_a,\gamma(1)=q_b}
  \max_{s\in[0,1]}\bar E(\gamma(s)).
\end{equation}
找到最大能量不超过 $h$ 的完整路径，可以证明这一对同分量；路径搜索失败不能证明
断开，声称断开需要 barrier 下界或全局拓扑证书。

去 gauge 切空间中的局部弱方向数可写为
\begin{equation}
  k_{\mathrm{weak}}(\tau)=\#\{j:\lambda_j\le\tau\}.
\end{equation}
它不等于全局连续解维数。还需 perturb--refit、continuation 和邻域秩稳定性检查。

\section{最小反例与证据不对称}

\begin{enumerate}
  \item 两条 raw 轨迹远离但对齐后相同：只是 gauge copies；
  \item 连通香蕉形集合只在两端被采到：有限聚类虚构了两个 component；
  \item 求解器有两个 basin，但存在合法低能路径：basin 不等于断开；
  \item 双端有效而均值无效：是速度多模态候选证据，仍需排除错误配对和标签噪声；
  \item 一个端点有效、一个无效：这是 selector problem；
  \item 双端和均值都有效：是连续可融合冗余，不必因“多个输出”强上 latent。
\end{enumerate}

\section*{结语}

理论的作用不是把模型复杂化，而是让模型复杂度与证据相配。对当前 VGGT 长短上下文
修复，最直接的顺序是：先把窗口输出变成可比较候选，再检验端点、方向和插值；由结果
决定确定性融合、selector 或 V-RFM。物理解空间研究继续提供更深的几何边界，但不再
被误写成 V-RFM 的唯一入场券。

\bibliographystyle{plainnat}
\bibliography{doc/references/camera_refiner_references}

\end{document}
